This chapter is freely inspired by the work of Bonneau, Faraut and Valent~\cite{Bonneau_2001}.
In physics, self-adjoint operators are
fundamental because they ensure that the associated observables have real
eigenvalues, and real eigenvalues are essential for physical interpretations of
measurements.

Most of our experience and intuition is predicated on the self-adjointness of the operators, and when this fails, the intelligibility of the theory goes with it.

Not all operators in quantum mechanics are self-adjoint by default, it may
depend on the operator itself or its considered domain.\footnote{For example,
    the momentum operator in one dimension is not self-adjoint on its maximum
    domain because its domain is not the entire Hilbert space, \textit{i.e. it's unbounded}.}

The method of self-adjoint extension provides a systematic approach to extend
or restrict the domain of a non-self-adjoint operator to make it self-adjoint.

The process we are about to describe can be summarized as follows: we begin with a symmetric operator $(A,\mathcal{D}(A))$ and gradually extend $\mathcal{D}(A)$ while simultaneously narrowing down $\mathcal{D}({A^\dagger})$ until the two domains become identical. This procedure gives rise to a set of important questions, which were initially tackled by Weyl and subsequently extended by von Neumann and Stone: 

How can one determine if a given operator allows for a self-adjoint extension? 

Is the extension, if it exists, unique? 

How is the self-adjoint domain constructed?


\vspace{4mm}
\noindent
This \textit{Mathematical Tool} chapter aims to answer those questions, but before we can delve into the essence of the method, we need some definitions.

\paragraph{Definition} Let us consider a Hilbert space $\mathcal{H}$. An operator $(A, \mathcal{D}(A))$ defined
on $\mathcal{H}$ is said to be densely defined if the subset $\mathcal{D}(A)$ is dense in
$\mathcal{H}$, meaning that for any $\psi \in \mathcal{H}$, one can find a sequence $\varphi_n$ in
$\mathcal{D}(A)$ that converges in norm to $\psi$.

\paragraph{Definition} An operator $(A, \mathcal{D}(A))$ is said to be closed if there exists a
sequence $\varphi_n$ in $\mathcal{D}(A)$ such that
\[
    \lim_{n \to \infty} \varphi_n = \varphi, \quad \lim_{n \to \infty} A\varphi_n = \psi,
\]
then it follows that $\varphi \in \mathcal{D}(A)$ and $A\varphi = \psi$.

Fortunately all the operators that we use in the rest of the essay are dense and closed, so we shouldn't worry about those conditions.

\paragraph{Definition} The adjoint operator of a (generally unbounded) operator $H$ with a dense
domain $\mathcal{D}(H)$ is defined as follows. The domain
$\mathcal{D}(H^\dagger)$ is the space of functions $\psi$ such that the linear
form
\[
    \varphi \mapsto \langle \psi| H\varphi \rangle
\]
is continuous with respect to the norm of $H$. This implies the existence of
$\psi^\dagger \in H$ such that
\[
    \langle\psi| H\varphi \rangle = \langle\psi^\dagger| \varphi \rangle .
\]
Hence, $H^\dagger\psi = \psi^\dagger$. Notably, the adjoint of any densely
defined operator is closed.

\paragraph{Definition} An operator $(H, \mathcal{D}(H))$ is said to be symmetric if for all $\varphi,
    \psi \in \mathcal{D}(H)$, we have
\[
    \langle H\varphi| \psi \rangle = \langle \varphi| H\psi \rangle.
\]
If $\mathcal{D}(H)$ is dense, this is equivalent to saying that $(H^\dagger,
    \mathcal{D}(H))$ is an extension of $(H, \mathcal{D}(H))$.

\paragraph{Definition} The operator $H$ with dense domain $\mathcal{D}(H)$ is considered self-adjoint
if $\mathcal{D}(H^\dagger) = \mathcal{D}(H)$ and $H^\dagger = H$.

\clearpage

\subsection{The Deficiency Indices}

In this section, we assume that $(A, \mathcal{D}(A))$ is densely defined,
symmetric, and closed, and let $(A^\dagger, \mathcal{D}(A^\dagger))$ be its
adjoint.

We define the deficiency subspaces $\mathcal{N}^+$ and $\mathcal{N}^-$ as
follows:
\[
    \mathcal{N}^+ = \{\psi \in \mathcal{D}(A^\dagger), \quad A^\dagger\psi = z^+\psi, \quad \text{ Im } z^+ > 0\},
\]
\[
    \mathcal{N}^- = \{\psi \in \mathcal{D}(A^\dagger), \quad A^\dagger\psi = z^-\psi, \quad \text{ Im } z^- < 0\},
\]

with respective dimensions $n^+$ and $n^-$. These are referred to as the
deficiency indices of the operator $A$ and will be denoted by the ordered pair
$(n^+, n^-)$.

Importantly, $n^+$ (or $n^-$) is independent of the choice of $z^+$ (or $z^-$)
as long as it lies in the upper (or lower) half-plane. Consequently, one simple
way to determine $(n^+, n^-)$ is to take $z^+ = i\eta$ and $z^- = -i\eta$
with an arbitrary strictly positive constant $\eta$ for dimensional reasons.

\vspace{4mm}
\noindent
The following theorem, by Weyl and von Neumann, is of primary importance.

\subsubsection*{Theorem:}

For an operator $A$ with deficiency indices $(n^+, n^-)$, there are three possibilities:
\begin{enumerate}
    \item If $n^+ = n^- = 0$, then $A$ is self-adjoint (in fact, this is a necessary and
          sufficient condition).
    \item If $n^+ = n^- = n \geq 1$, then $A$ has infinitely many self-adjoint
          extensions, parametrized by a unitary $n \times n$ matrix (\textit{i.e. $n^2$ real parameters}).
    \item If $n^+ \neq n^-$, then $A$ has no self-adjoint extension.
\end{enumerate}

\noindent
The application of this theorem to differential operators requires substantial
work. Even if we start with an operator $P$ that is formally self-adjoint, this
does not guarantee that $P$ is truly self-adjoint because the domains
$\mathcal{D}(P)$ and $\mathcal{D}(P^\dagger)$ will generally be different.

For a given differential operator $P$, three problems need to be solved:
\begin{enumerate}
    \item Find a domain $\mathcal{D}(P)$ for which the formally self-adjoint operator $P$
          is symmetric and closed.
    \item Compute its adjoint $(P^\dagger, \mathcal{D}(P^\dagger))$ and determine the
          deficiency indices of $P^\dagger$.
    \item When they do exist, describe the domains of all the self-adjoint extensions.
\end{enumerate}

\vspace{4mm}
\noindent
In this section, we have briefly described the process to understand the existence, the uniqueness and the dimension of the space of self-adjoint extensions.

\subsection{The Boundary Conditions}

Now we need to provide a method to compute boundary conditions that make our operator $H$ self-adjoint.

For simplicity we assume that the deficiency indices are $n_+ = n_- = 1$, \textit{i.e. $(1,1)$}. Then we expect to find $\phi_+$ and $\phi_-$ as eigenfunction with respect to the two imaginary eigenvalues, \textit{i.e. two generators of the subspaces $ \mathcal{N}_\pm $}.

We also know that the unitary matrix that relates the two one dimensional subspaces $ \mathcal{N}_\pm $ is one-dimensional, \textit{i.e. a complex number of modulus one}. Hence, the self-adjoint extensions depend on one parameter only.

We set $U = \lambda$ with $|\lambda| = 1$, and define $\Phi_\lambda = \phi_+ + \lambda \phi_-$. The prescription for obtaining the boundary conditions is simply to require that:
\begin{align} 
    \label{eq:boundary_condition}
    \langle H \Phi_\lambda | \psi \rangle = \langle \Phi_\lambda | H \psi \rangle
\end{align}

\noindent
All and only the $\psi$ that satisfy this condition are in the domain of the $\lambda$-self-adjoint extension of the operator $H$, \textit{i.e. $(H,\mathcal{D}_\lambda(H))$}.